% ══════════════════════════════════════════════════════════════════════════════
% Chapter 3: Methodology
% ══════════════════════════════════════════════════════════════════════════════

This chapter presents the detailed methodology of the PISL framework as proposed by Tao~\textit{et~al.}~\cite{tao2024unsupervised}. The overall architecture is illustrated in Figure~\ref{fig:overview}.

\begin{figure}[H]
    \centering
    \includegraphics[width=\textwidth]{overview.png}
    \caption{Overview of the PISL framework~\cite{tao2024unsupervised}. The framework consists of a ResNet-50 backbone for feature extraction, a Spatial Diffusion Model (SDM) that generates transformation parameters through a denoising process, and a sampler that extracts semantically meaningful patches. Training is guided by five loss functions: cross-entropy with label smoothing ($\mathcal{L}_{ces}$), soft triplet loss ($\mathcal{L}_{tri}$), camera-aware contrastive loss ($\mathcal{L}_{cam}$), semantic controlled diffusion loss ($\mathcal{L}_{scd}$), and patch semantic consistency loss ($\mathcal{L}_{psc}$).}
    \label{fig:overview}
\end{figure}

\section{Overall Framework}

Given an unlabelled dataset $\mathcal{X} = \{x_i\}_{i=1}^{N}$ with camera annotations $\{c_i\}_{i=1}^{N}$ but without identity labels, the PISL framework iteratively performs four steps:

\begin{enumerate}
    \item \textbf{Feature Extraction:} Extract global features $\mathbf{f}_g \in \mathbb{R}^{D}$ and part features $\mathbf{f}_p \in \mathbb{R}^{D \times P}$ using the backbone and Spatial Diffusion Model, where $D$ is the feature dimension and $P$ is the number of parts.
    \item \textbf{Pseudo Label Generation:} Cluster global features using DBSCAN~\cite{ester1996dbscan} with Jaccard distance to assign pseudo identity labels.
    \item \textbf{Semantic Consistency Scoring:} Compute consistency scores between global and part features using Maximum Mean Discrepancy (MMD)~\cite{gretton2012mmd}.
    \item \textbf{Model Training:} Update the model using a combination of classification, metric learning, camera-aware contrastive, and diffusion-based losses.
\end{enumerate}

\section{Backbone and Feature Extraction}

The backbone network is a ResNet-50~\cite{he2016deep} pretrained on ImageNet~\cite{deng2009imagenet}. Given an input image $x$, the backbone extracts a feature map $\mathbf{F} \in \mathbb{R}^{2048 \times H \times W}$ from the last convolutional layer, where $H$ and $W$ are the spatial dimensions. A global average pooling operation produces the global feature vector:
\begin{equation}
    \mathbf{f}_g = \text{GAP}(\mathbf{F}) \in \mathbb{R}^{2048}
\end{equation}
which is subsequently normalized via batch normalization and L2 normalization for distance computation.

\section{Spatial Diffusion Model (SDM)}
\label{sec:sdm}

The Spatial Diffusion Model is the central innovation of PISL. Instead of generating patches through a learned static spatial transformer, PISL introduces a \textit{diffusion process} over the spatial transformer parameters $\boldsymbol{\theta}$.

\subsection{Localization Network}

A localization network processes the backbone feature map $\mathbf{F}$ to produce initial transformation parameters $\boldsymbol{\theta}_0 \in \mathbb{R}^{P \times 2 \times 3}$, where $P$ is the number of parts (set to 3 in practice). The localization network consists of:
\begin{itemize}
    \item A convolutional layer: Conv2d($2048 \rightarrow 4096$, kernel $3 \times 3$)
    \item Batch normalization and ReLU activation
    \item Max pooling (kernel $3 \times 3$, stride 2)
    \item Adaptive average pooling to $1 \times 1$
    \item Two fully connected layers: FC($4096 \rightarrow 512$) $\rightarrow$ FC($512 \rightarrow 3 \times 2 \times 3$)
\end{itemize}

Each $\boldsymbol{\theta}_0^{(i)}$ defines an affine transformation for the $i$-th part:
\begin{equation}
    \boldsymbol{\theta}_0^{(i)} = \begin{pmatrix} s_x & 0 & t_x \\ 0 & s_y & t_y \end{pmatrix}
\end{equation}
where $s_x, s_y$ control the scale and $t_x, t_y$ control the translation of each patch. The parameters are initialized to produce three vertically stacked horizontal strips (head, torso, legs).

\subsection{Forward Diffusion Process}

The forward process gradually adds Gaussian noise to the initial parameters $\boldsymbol{\theta}_0$ over $T = 1000$ timesteps according to a variance schedule $\{\beta_t\}_{t=1}^{T}$:
\begin{equation}
    q(\boldsymbol{\theta}_t | \boldsymbol{\theta}_{t-1}) = \mathcal{N}\left(\boldsymbol{\theta}_t;\ \sqrt{1-\beta_t}\,\boldsymbol{\theta}_{t-1},\ \beta_t\,\mathbf{I}\right)
\end{equation}

Using the reparameterization trick with $\alpha_t = 1 - \beta_t$ and $\bar{\alpha}_t = \prod_{s=1}^{t}\alpha_s$, the noisy parameters at any timestep $t$ can be directly computed as:
\begin{equation}
    \boldsymbol{\theta}_t = \sqrt{\bar{\alpha}_t}\,\boldsymbol{\theta}_0 + \sqrt{1 - \bar{\alpha}_t}\,\boldsymbol{\epsilon}, \quad \boldsymbol{\epsilon} \sim \mathcal{N}(\mathbf{0}, \mathbf{I})
\end{equation}

A \textit{cosine} $\beta$-schedule~\cite{ho2020denoising} is employed with offset $s = 0.008$:
\begin{equation}
    \bar{\alpha}_t = \frac{f(t)}{f(0)}, \quad f(t) = \cos\left(\frac{t/T + s}{1 + s} \cdot \frac{\pi}{2}\right)^2
\end{equation}

This schedule provides a more gradual noise addition compared to the linear schedule, preserving signal information for longer during the forward process.

\subsection{Reverse Denoising Process}

A denoising network $\boldsymbol{\Phi}$ is trained to predict the noise $\boldsymbol{\epsilon}$ added at each timestep, enabling recovery of the clean parameters:
\begin{equation}
    \hat{\boldsymbol{\theta}}_0 = \frac{1}{\sqrt{\bar{\alpha}_t}} \left( \boldsymbol{\theta}_t - \sqrt{1 - \bar{\alpha}_t} \cdot \boldsymbol{\Phi}(\boldsymbol{\theta}_t, t) \right)
\end{equation}

During inference, the denoising process starts from pure Gaussian noise $\boldsymbol{\theta}_T \sim \mathcal{N}(\mathbf{0}, \mathbf{I})$ and iteratively denoises to obtain the final semantic parameters $\hat{\boldsymbol{\theta}}_0$.

\subsection{Patch Sampling via Spatial Transformer}

The denoised parameters $\hat{\boldsymbol{\theta}}_0$ are used to construct a sampling grid via the Spatial Transformer Network~\cite{jaderberg2015spatial}:
\begin{equation}
    \mathbf{G}_i = \hat{\boldsymbol{\theta}}_0^{(i)} \cdot \mathbf{G}_{\text{base}}
\end{equation}
where $\mathbf{G}_{\text{base}}$ is the regular grid and $\mathbf{G}_i$ is the transformed grid for part $i$. Bilinear sampling extracts the part feature map $\mathbf{F}_i$ from the backbone feature map $\mathbf{F}$. Each part feature is then processed through global average pooling and batch normalization to produce the part feature vector $\mathbf{f}_p^{(i)} \in \mathbb{R}^{D}$.

\section{Semantic Controlled Diffusion (SCD) Loss}

The SCD loss guides the denoising process toward semantically meaningful transformations. It consists of three components:

\subsection{Noise Prediction Loss}

Trains the denoising network to accurately predict the noise added at each timestep:
\begin{equation}
    \mathcal{L}_{\text{noise}} = \|\boldsymbol{\Phi}(\boldsymbol{\theta}_t, t) - \boldsymbol{\epsilon}\|_2^2
\end{equation}

\subsection{Theta Consistency Loss}

Ensures the denoised parameters remain close to the localization network's output, combining L2 distance and cosine similarity:
\begin{equation}
    \mathcal{L}_{\text{consist}} = \|\hat{\boldsymbol{\theta}}_0 - \boldsymbol{\theta}_0\|_2^2 + 0.1 \cdot \left(1 - \cos(\hat{\boldsymbol{\theta}}_0, \boldsymbol{\theta}_0)\right)
\end{equation}

\subsection{Part Contrastive Loss}

Encourages discriminative part features through contrastive learning with learned part weights:
\begin{equation}
    \mathcal{L}_{\text{part}} = -\sum_{i=1}^{P} w_i \cdot \frac{1}{|\mathcal{P}_i|} \sum_{j \in \mathcal{P}_i} \log \frac{\exp(\mathbf{f}_p^{(i)} \cdot \mathbf{f}_{p,j}^{(i)} / \tau)}{\sum_k \exp(\mathbf{f}_p^{(i)} \cdot \mathbf{f}_{p,k}^{(i)} / \tau)}
\end{equation}
where $w_i = \exp(-\mathcal{L}_{\text{part}}^{(i)}) / \sum_j \exp(-\mathcal{L}_{\text{part}}^{(j)})$ is a learned weight based on part discriminability, $\mathcal{P}_i$ denotes positive samples, and $\tau = 0.07$ is the temperature parameter.

The combined SCD loss is:
\begin{equation}
    \mathcal{L}_{scd} = \mathcal{L}_{\text{noise}} + \lambda_c \cdot \mathcal{L}_{\text{consist}} + \lambda_p \cdot \mathcal{L}_{\text{part}}
\end{equation}
where $\lambda_c = 0.5$ and $\lambda_p = 0.3$ are the weighting hyperparameters.

\section{Patch Semantic Consistency (PSC) Loss}

The PSC loss leverages the semantic consistency between global and part features to refine pseudo labels.

\subsection{Consistency Measurement via MMD}

The consistency between global and part features is measured using Maximum Mean Discrepancy (MMD)~\cite{gretton2012mmd}. For each sample $j$ and part $i$, the $k$-nearest neighbour sets are computed in both the global feature space ($\mathcal{G}_j$) and the $i$-th part feature space ($\mathcal{P}_j^{(i)}$). The consistency score is:
\begin{equation}
    s_j^{(i)} = \exp\left(-\text{MMD}^2(\mathcal{G}_j, \mathcal{P}_j^{(i)})\right)
\end{equation}
where $\text{MMD}^2(\mathcal{G}_j, \mathcal{P}_j^{(i)}) = \|\mu_{\mathcal{G}_j} - \mu_{\mathcal{P}_j^{(i)}}\|_{\mathcal{H}}^2$ is the squared MMD in a reproducing kernel Hilbert space.

A high consistency score indicates that the global and part features agree on the neighbourhood structure, suggesting more reliable pseudo labels.

\subsection{PSC Label Refinement (PSC-LR)}

Uses the weighted ensemble of part predictions to refine global pseudo labels:
\begin{equation}
    \tilde{y} = \lambda_{\text{ref}} \cdot y + (1 - \lambda_{\text{ref}}) \cdot \sum_{i=1}^{P} w_i \cdot p_i
\end{equation}
where $y$ is the one-hot pseudo label, $p_i$ is the softmax prediction from part $i$, $w_i$ is the consistency-based weight obtained by normalizing $s_j^{(i)}$ via softmax, and $\lambda_{\text{ref}} = 0.5$ controls the refinement strength.

\subsection{PSC Label Smoothing (PSC-LS)}

Applies consistency-aware label smoothing to part classifiers:
\begin{equation}
    \mathcal{L}_{\text{psc-ls}}^{(i)} = s_j^{(i)} \cdot \mathcal{L}_{ce}(\mathbf{f}_p^{(i)}, y) + (1 - s_j^{(i)}) \cdot D_{KL}(p_i \,\|\, u)
\end{equation}
where $u$ is the uniform distribution. When the consistency score is high, the loss emphasizes the cross-entropy term; when it is low, the KL divergence toward the uniform distribution acts as a regularizer, preventing the model from overfitting to noisy labels.

\section{Camera-Aware Contrastive Learning}

To address the domain gap introduced by different camera conditions, the framework incorporates camera-aware contrastive learning~\cite{wang2021cap}. Camera-level proxies are constructed for each (pseudo identity, camera) pair. The camera contrastive loss is:
\begin{equation}
    \mathcal{L}_{cam} = -\sum_{i=1}^{B} \frac{1}{|\mathcal{P}_i^+|} \sum_{j \in \mathcal{P}_i^+} \log \frac{\exp(s_{ij}/\tau)}{\sum_{k \in \mathcal{P}_i^+ \cup \mathcal{N}_i^-} \exp(s_{ik}/\tau)}
\end{equation}
where $\mathcal{P}_i^+$ contains positive proxies of the same identity but \textit{different} cameras, $\mathcal{N}_i^-$ contains the top-$K$ hard negative proxies (set to $K=50$), and $\tau = 0.07$ is the temperature. This loss explicitly encourages the model to learn features that are invariant to camera-specific biases such as illumination, background, and resolution differences.

\section{Total Training Objective}

The overall training objective combines all loss components:
\begin{equation}
    \mathcal{L} = \mathcal{L}_{ces} + \mathcal{L}_{tri} + W_{cam} \cdot \mathcal{L}_{cam} + W_{diff} \cdot \mathcal{L}_{scd} + \mathcal{L}_{psc}
\end{equation}
where:
\begin{itemize}
    \item $\mathcal{L}_{ces}$: Cross-entropy with label smoothing~\cite{szegedy2016rethinking} for global feature classification.
    \item $\mathcal{L}_{tri}$: Soft triplet loss~\cite{hermans2017triplet} for metric learning.
    \item $W_{cam} = 0.5$: Weight for the camera-aware contrastive loss.
    \item $W_{diff} = 0.1$: Weight for the semantic controlled diffusion loss.
    \item $\mathcal{L}_{psc}$: Combined PSC-LR and PSC-LS losses.
\end{itemize}
